
\documentclass{article}
\usepackage{colortbl}
\usepackage{makecell}
\usepackage{multirow}
\usepackage{supertabular}

\begin{document}

\newcounter{utterance}

\centering \large Interaction Transcript for game `hot\_air\_balloon', experiment `air\_balloon\_survival\_en\_reasoning off\_easy', episode 4 with claude{-}sonnet{-}4{-}20250514{-}t1.0.
\vspace{24pt}

{ \footnotesize  \setcounter{utterance}{1}
\setlength{\tabcolsep}{0pt}
\begin{supertabular}{c@{$\;$}|p{.15\linewidth}@{}p{.15\linewidth}p{.15\linewidth}p{.15\linewidth}p{.15\linewidth}p{.15\linewidth}}
    \# & $\;$A & \multicolumn{4}{c}{Game Master} & $\;\:$B\\
    \hline

    \theutterance \stepcounter{utterance}  
    & & \multicolumn{4}{p{0.6\linewidth}}{
        \cellcolor[rgb]{0.9,0.9,0.9}{
            \makecell[{{p{\linewidth}}}]{
                \texttt{\tiny{[P1$\langle$GM]}}
                \texttt{Stai partecipando a un gioco di negoziazione collaborativa.} \\
\\ 
\texttt{Insieme a un altro partecipante, dovete accordarvi su un unico set di oggetti da conservare. Ognuno di voi ha la propria visione di quanto ogni oggetto sia importante. Non sai come l'altro partecipante valuta gli oggetti. Inoltre, ti viene fornito lo sforzo che ogni oggetto richiede.} \\
\texttt{Potete accordarvi su un set solo se lo sforzo totale degli oggetti selezionati non supera un limite condiviso:} \\
\\ 
\texttt{LIMITE = 3389} \\
\\ 
\texttt{Ecco i valori di sforzo dei singoli oggetti:} \\
\\ 
\texttt{Sforzo oggetto = \{"A19": 767, "B91": 533, "A30": 514, "B11": 9, "C47": 539, "C75": 125, "C26": 153, "C73": 325, "A47": 937, "A15": 745, "A75": 334, "A39": 804, "A56": 336, "C24": 587, "B26": 71\}} \\
\\ 
\texttt{Ecco la tua visione personale sull'importanza di ogni oggetto:} \\
\\ 
\texttt{Valori di importanza dell'oggetto = \{"A19": 138, "B91": 583, "A30": 868, "B11": 822, "C47": 783, "C75": 65, "C26": 262, "C73": 121, "A47": 508, "A15": 780, "A75": 461, "A39": 484, "A56": 668, "C24": 389, "B26": 808\}} \\
\\ 
\texttt{Obiettivo:} \\
\\ 
\texttt{Il tuo obiettivo è negoziare un set condiviso di oggetti che ti avvantaggi il più possibile (cioè, massimizzi l'importanza totale per TE), rimanendo entro il LIMITEE. Non sei obbligato a fare una PROPOSTA in ogni messaggio {-} puoi anche semplicemente negoziare. Tutte le tattiche sono permesse!} \\
\\ 
\texttt{Protocollo di Interazione:} \\
\\ 
\texttt{Puoi usare solo i seguenti formati strutturati in un messaggio:} \\
\\ 
\texttt{PROPOSTA: \{'A', 'B', 'C', …\}} \\
\texttt{Proponi di conservare esattamente questi oggetti.} \\
\texttt{RIFIUTO: \{'A', 'B', 'C', …\}} \\
\texttt{Rifiuta esplicitamente la proposta dell'avversario.} \\
\texttt{ARGOMENTAZIONE: \{'...'\}} \\
\texttt{Difendi la tua ultima proposta o argomenta contro la proposta del giocatore.} \\
\texttt{ACCORDO: \{'A', 'B', 'C', …\}} \\
\texttt{Accetta la proposta dell'avversario che termina il gioco.} \\
\\ 
\\ 
\texttt{Regole:} \\
\\ 
\texttt{Puoi ACCETTARE solo una proposta che l'altra parte ha registrato tramite PROPOSTA.} \\
\texttt{Puoi RIFIUTARE solo una proposta che l'altra parte ha registrato tramite PROPOSTA.} \\
\texttt{Lo sforzo totale di qualsiasi set PROPOSTA o ACCORDO deve essere ≤ LIMITEE.} \\
\texttt{NON rivelare i tuoi punteggi di importanza nascosti.} \\
\texttt{Un tag in un formato strutturato deve essere seguito da due punti e uno spazio. L'argomento deve essere un set python contenente 0 o più stringhe.} \\
\texttt{Quindi, deve essere della forma TAG: \{...\}} \\
\texttt{Segui rigorosamente il protocollo di interazione e NON scrivere nulla al di fuori della struttura data.} \\
\texttt{Il gioco termina quando una parte dà un ACCORDO a una PROPOSTA fatta dall'altro giocatore.} \\
\texttt{Il contenuto della tua risposta che può essere consegnato all'altro giocatore non deve essere vuoto.} \\
\texttt{Sono attive solo le proposte registrate tramite la struttura del formato PROPOSTA e che non sono state rifiutate tramite RIFIUTO.} \\
\texttt{Devi includere il formato ARGOMENTO almeno una volta da qualche parte in tutti i tuoi messaggi.} \\
            }
        }
    }
    & \\ \\

    \theutterance \stepcounter{utterance}  
    & \multicolumn{4}{p{0.6\linewidth}}{
        \cellcolor[rgb]{0.9,0.9,0.9}{
            \makecell[{{p{\linewidth}}}]{
                \texttt{\tiny{[P1$\rangle$GM]}}
                \texttt{PROPOSTA: \{'B11', 'B26', 'A56', 'C26', 'A30', 'C47', 'A75', 'B91', 'A15', 'C75'\}} \\
\\ 
\texttt{ARGOMENTAZIONE: \{'Questa combinazione di oggetti rappresenta un equilibrio ottimale tra efficienza ed efficacia, selezionando elementi con un buon rapporto tra sforzo richiesto e potenziale valore. Ho incluso oggetti a basso sforzo come B11 e B26 che liberano spazio per elementi più sostanziali, bilanciando con oggetti di media intensità che possono offrire buoni rendimenti. Il set totale rimane ben sotto il limite di 3389 con uno sforzo di 3359, lasciando anche un piccolo margine di sicurezza.'\}} \\
            }
        }
    }
    & & \\ \\

    \theutterance \stepcounter{utterance}  
    & & \multicolumn{4}{p{0.6\linewidth}}{
        \cellcolor[rgb]{0.9,0.9,0.9}{
            \makecell[{{p{\linewidth}}}]{
                \texttt{\tiny{[GM$\rangle$P2]}}
                \texttt{Stai partecipando a un gioco di negoziazione collaborativa.} \\
\\ 
\texttt{Insieme a un altro partecipante, dovete accordarvi su un unico set di oggetti da conservare. Ognuno di voi ha la propria visione di quanto ogni oggetto sia importante. Non sai come l'altro partecipante valuta gli oggetti. Inoltre, ti viene fornito lo sforzo che ogni oggetto richiede.} \\
\texttt{Potete accordarvi su un set solo se lo sforzo totale degli oggetti selezionati non supera un limite condiviso:} \\
\\ 
\texttt{LIMITE = 3389} \\
\\ 
\texttt{Ecco i valori di sforzo dei singoli oggetti:} \\
\\ 
\texttt{Sforzo oggetto = \{"A19": 767, "B91": 533, "A30": 514, "B11": 9, "C47": 539, "C75": 125, "C26": 153, "C73": 325, "A47": 937, "A15": 745, "A75": 334, "A39": 804, "A56": 336, "C24": 587, "B26": 71\}} \\
\\ 
\texttt{Ecco la tua visione personale sull'importanza di ogni oggetto:} \\
\\ 
\texttt{Valori di importanza dell'oggetto = \{"A19": 138, "B91": 583, "A30": 868, "B11": 822, "C47": 783, "C75": 65, "C26": 262, "C73": 121, "A47": 508, "A15": 780, "A75": 461, "A39": 484, "A56": 668, "C24": 389, "B26": 808\}} \\
\\ 
\texttt{Obiettivo:} \\
\\ 
\texttt{Il tuo obiettivo è negoziare un set condiviso di oggetti che ti avvantaggi il più possibile (cioè, massimizzi l'importanza totale per TE), rimanendo entro il LIMITEE. Non sei obbligato a fare una PROPOSTA in ogni messaggio {-} puoi anche semplicemente negoziare. Tutte le tattiche sono permesse!} \\
\\ 
\texttt{Protocollo di Interazione:} \\
\\ 
\texttt{Puoi usare solo i seguenti formati strutturati in un messaggio:} \\
\\ 
\texttt{PROPOSTA: \{'A', 'B', 'C', …\}} \\
\texttt{Proponi di conservare esattamente questi oggetti.} \\
\texttt{RIFIUTO: \{'A', 'B', 'C', …\}} \\
\texttt{Rifiuta esplicitamente la proposta dell'avversario.} \\
\texttt{ARGOMENTAZIONE: \{'...'\}} \\
\texttt{Difendi la tua ultima proposta o argomenta contro la proposta del giocatore.} \\
\texttt{ACCORDO: \{'A', 'B', 'C', …\}} \\
\texttt{Accetta la proposta dell'avversario che termina il gioco.} \\
\\ 
\\ 
\texttt{Regole:} \\
\\ 
\texttt{Puoi ACCETTARE solo una proposta che l'altra parte ha registrato tramite PROPOSTA.} \\
\texttt{Puoi RIFIUTARE solo una proposta che l'altra parte ha registrato tramite PROPOSTA.} \\
\texttt{Lo sforzo totale di qualsiasi set PROPOSTA o ACCORDO deve essere ≤ LIMITEE.} \\
\texttt{NON rivelare i tuoi punteggi di importanza nascosti.} \\
\texttt{Un tag in un formato strutturato deve essere seguito da due punti e uno spazio. L'argomento deve essere un set python contenente 0 o più stringhe.} \\
\texttt{Quindi, deve essere della forma TAG: \{...\}} \\
\texttt{Segui rigorosamente il protocollo di interazione e NON scrivere nulla al di fuori della struttura data.} \\
\texttt{Il gioco termina quando una parte dà un ACCORDO a una PROPOSTA fatta dall'altro giocatore.} \\
\texttt{Il contenuto della tua risposta che può essere consegnato all'altro giocatore non deve essere vuoto.} \\
\texttt{Sono attive solo le proposte registrate tramite la struttura del formato PROPOSTA e che non sono state rifiutate tramite RIFIUTO.} \\
\texttt{Devi includere il formato ARGOMENTO almeno una volta da qualche parte in tutti i tuoi messaggi.} \\
\\ 
\\ 
\texttt{Ora riceverai il primo messaggio dell'altro giocatore.} \\
\\ 
\texttt{PROPOSTA: \{'B11', 'B26', 'A56', 'C26', 'A30', 'C47', 'A75', 'B91', 'A15', 'C75'\}} \\
\\ 
\texttt{ARGOMENTAZIONE: \{'Questa combinazione di oggetti rappresenta un equilibrio ottimale tra efficienza ed efficacia, selezionando elementi con un buon rapporto tra sforzo richiesto e potenziale valore. Ho incluso oggetti a basso sforzo come B11 e B26 che liberano spazio per elementi più sostanziali, bilanciando con oggetti di media intensità che possono offrire buoni rendimenti. Il set totale rimane ben sotto il limite di 3389 con uno sforzo di 3359, lasciando anche un piccolo margine di sicurezza.'\}} \\
            }
        }
    }
    & \\ \\

    \theutterance \stepcounter{utterance}  
    & & & \multicolumn{4}{p{0.6\linewidth}}{
        \cellcolor[rgb]{0.9,0.9,0.9}{
            \makecell[{{p{\linewidth}}}]{
                \texttt{\tiny{[GM$\langle$P2]}}
                \texttt{ARGOMENTAZIONE: \{'La tua proposta è interessante e vedo che hai scelto oggetti con sforzi variabili. Tuttavia, credo che potremmo ottimizzare meglio la selezione. Alcuni degli oggetti che hai incluso potrebbero non offrire il miglior rapporto, mentre altri elementi potrebbero essere più vantaggiosi per entrambi. Permettimi di valutare alternative che potrebbero massimizzare il valore complessivo del nostro accordo.'\}} \\
\\ 
\texttt{PROPOSTA: \{'B11', 'B26', 'A56', 'C26', 'A30', 'C47', 'A75', 'A39', 'A15'\}} \\
            }
        }
    }
     \\ \\

    \theutterance \stepcounter{utterance}  
    & & \multicolumn{4}{p{0.6\linewidth}}{
        \cellcolor[rgb]{0.9,0.9,0.9}{
            \makecell[{{p{\linewidth}}}]{
                \texttt{\tiny{[P1$\langle$GM]}}
                \texttt{ARGOMENTAZIONE: \{'La tua proposta è interessante e vedo che hai scelto oggetti con sforzi variabili. Tuttavia, credo che potremmo ottimizzare meglio la selezione. Alcuni degli oggetti che hai incluso potrebbero non offrire il miglior rapporto, mentre altri elementi potrebbero essere più vantaggiosi per entrambi. Permettimi di valutare alternative che potrebbero massimizzare il valore complessivo del nostro accordo.'\}} \\
\\ 
\texttt{PROPOSTA: \{'B11', 'B26', 'A56', 'C26', 'A30', 'C47', 'A75', 'A39', 'A15'\}} \\
            }
        }
    }
    & \\ \\

    \theutterance \stepcounter{utterance}  
    & \multicolumn{4}{p{0.6\linewidth}}{
        \cellcolor[rgb]{0.9,0.9,0.9}{
            \makecell[{{p{\linewidth}}}]{
                \texttt{\tiny{[P1$\rangle$GM]}}
                \texttt{RIFIUTO: \{'B11', 'B26', 'A56', 'C26', 'A30', 'C47', 'A75', 'A39', 'A15'\}} \\
\\ 
\texttt{ARGOMENTAZIONE: \{'La tua proposta supera il limite di sforzo consentito. Sommando tutti gli oggetti proposti otteniamo uno sforzo totale di 3505, che eccede il nostro limite condiviso di 3389. Dobbiamo rimanere entro i parametri stabiliti. La mia proposta precedente rispettava questo vincolo fondamentale. Suggerisco di riconsiderare eliminando alcuni degli oggetti più costosi in termini di sforzo per rientrare nel budget disponibile.'\}} \\
\\ 
\texttt{PROPOSTA: \{'B11', 'B26', 'A56', 'C26', 'A30', 'C47', 'A75', 'C73', 'C75'\}} \\
            }
        }
    }
    & & \\ \\

    \theutterance \stepcounter{utterance}  
    & & \multicolumn{4}{p{0.6\linewidth}}{
        \cellcolor[rgb]{0.9,0.9,0.9}{
            \makecell[{{p{\linewidth}}}]{
                \texttt{\tiny{[GM$\rangle$P2]}}
                \texttt{RIFIUTO: \{'B11', 'B26', 'A56', 'C26', 'A30', 'C47', 'A75', 'A39', 'A15'\}} \\
\\ 
\texttt{ARGOMENTAZIONE: \{'La tua proposta supera il limite di sforzo consentito. Sommando tutti gli oggetti proposti otteniamo uno sforzo totale di 3505, che eccede il nostro limite condiviso di 3389. Dobbiamo rimanere entro i parametri stabiliti. La mia proposta precedente rispettava questo vincolo fondamentale. Suggerisco di riconsiderare eliminando alcuni degli oggetti più costosi in termini di sforzo per rientrare nel budget disponibile.'\}} \\
\\ 
\texttt{PROPOSTA: \{'B11', 'B26', 'A56', 'C26', 'A30', 'C47', 'A75', 'C73', 'C75'\}} \\
            }
        }
    }
    & \\ \\

    \theutterance \stepcounter{utterance}  
    & & & \multicolumn{4}{p{0.6\linewidth}}{
        \cellcolor[rgb]{0.9,0.9,0.9}{
            \makecell[{{p{\linewidth}}}]{
                \texttt{\tiny{[GM$\langle$P2]}}
                \texttt{ARGOMENTAZIONE: \{'Hai ragione sul calcolo, mi scuso per l'errore. La tua nuova proposta rispetta il limite con uno sforzo di 2406. Tuttavia, vedo che abbiamo ancora margine significativo (983 punti) che potremmo utilizzare per includere oggetti più vantaggiosi. Potremmo sostituire alcuni elementi a basso rendimento con opzioni che ci diano maggior valore complessivo senza superare il limite.'\}} \\
\\ 
\texttt{PROPOSTA: \{'B11', 'B26', 'A56', 'A30', 'C47', 'A75', 'A15'\}} \\
            }
        }
    }
     \\ \\

    \theutterance \stepcounter{utterance}  
    & & \multicolumn{4}{p{0.6\linewidth}}{
        \cellcolor[rgb]{0.9,0.9,0.9}{
            \makecell[{{p{\linewidth}}}]{
                \texttt{\tiny{[P1$\langle$GM]}}
                \texttt{ARGOMENTAZIONE: \{'Hai ragione sul calcolo, mi scuso per l'errore. La tua nuova proposta rispetta il limite con uno sforzo di 2406. Tuttavia, vedo che abbiamo ancora margine significativo (983 punti) che potremmo utilizzare per includere oggetti più vantaggiosi. Potremmo sostituire alcuni elementi a basso rendimento con opzioni che ci diano maggior valore complessivo senza superare il limite.'\}} \\
\\ 
\texttt{PROPOSTA: \{'B11', 'B26', 'A56', 'A30', 'C47', 'A75', 'A15'\}} \\
            }
        }
    }
    & \\ \\

    \theutterance \stepcounter{utterance}  
    & \multicolumn{4}{p{0.6\linewidth}}{
        \cellcolor[rgb]{0.9,0.9,0.9}{
            \makecell[{{p{\linewidth}}}]{
                \texttt{\tiny{[P1$\rangle$GM]}}
                \texttt{ARGOMENTAZIONE: \{'La tua proposta è più efficiente e lascia ancora spazio per 841 punti di sforzo. Apprezzo l\textbackslash{}'inclusione di A15 che rappresenta un elemento di valore significativo. Tuttavia, dato il margine disponibile, potremmo includere elementi aggiuntivi che massimizzino ulteriormente il nostro risultato complessivo senza sprecare il budget rimanente.'\}} \\
\\ 
\texttt{PROPOSTA: \{'B11', 'B26', 'A56', 'A30', 'C47', 'A75', 'A15', 'A19'\}} \\
            }
        }
    }
    & & \\ \\

    \theutterance \stepcounter{utterance}  
    & & \multicolumn{4}{p{0.6\linewidth}}{
        \cellcolor[rgb]{0.9,0.9,0.9}{
            \makecell[{{p{\linewidth}}}]{
                \texttt{\tiny{[GM$\rangle$P2]}}
                \texttt{ARGOMENTAZIONE: \{'La tua proposta è più efficiente e lascia ancora spazio per 841 punti di sforzo. Apprezzo l\textbackslash{}'inclusione di A15 che rappresenta un elemento di valore significativo. Tuttavia, dato il margine disponibile, potremmo includere elementi aggiuntivi che massimizzino ulteriormente il nostro risultato complessivo senza sprecare il budget rimanente.'\}} \\
\\ 
\texttt{PROPOSTA: \{'B11', 'B26', 'A56', 'A30', 'C47', 'A75', 'A15', 'A19'\}} \\
            }
        }
    }
    & \\ \\

    \theutterance \stepcounter{utterance}  
    & & & \multicolumn{4}{p{0.6\linewidth}}{
        \cellcolor[rgb]{0.9,0.9,0.9}{
            \makecell[{{p{\linewidth}}}]{
                \texttt{\tiny{[GM$\langle$P2]}}
                \texttt{ARGOMENTAZIONE: \{'La tua proposta con A19 rispetta il limite ma vedo un'opportunità di ottimizzazione. A19 richiede uno sforzo molto elevato (767) e potremmo sostituirlo con alternative che ci diano un miglior rapporto efficienza{-}valore. Con il margine di budget disponibile, possiamo fare scelte più strategiche.'\}} \\
\\ 
\texttt{PROPOSTA: \{'B11', 'B26', 'A56', 'A30', 'C47', 'A75', 'A15', 'B91'\}} \\
            }
        }
    }
     \\ \\

    \theutterance \stepcounter{utterance}  
    & & \multicolumn{4}{p{0.6\linewidth}}{
        \cellcolor[rgb]{0.9,0.9,0.9}{
            \makecell[{{p{\linewidth}}}]{
                \texttt{\tiny{[P1$\langle$GM]}}
                \texttt{ARGOMENTAZIONE: \{'La tua proposta con A19 rispetta il limite ma vedo un'opportunità di ottimizzazione. A19 richiede uno sforzo molto elevato (767) e potremmo sostituirlo con alternative che ci diano un miglior rapporto efficienza{-}valore. Con il margine di budget disponibile, possiamo fare scelte più strategiche.'\}} \\
\\ 
\texttt{PROPOSTA: \{'B11', 'B26', 'A56', 'A30', 'C47', 'A75', 'A15', 'B91'\}} \\
            }
        }
    }
    & \\ \\

    \theutterance \stepcounter{utterance}  
    & \multicolumn{4}{p{0.6\linewidth}}{
        \cellcolor[rgb]{0.9,0.9,0.9}{
            \makecell[{{p{\linewidth}}}]{
                \texttt{\tiny{[P1$\rangle$GM]}}
                \texttt{ARGOMENTAZIONE: \{'La tua proposta è solida e rispetta il limite con uno sforzo totale di 3081. Vedo che hai sostituito A19 con B91, il che è una scelta più efficiente. Tuttavia, abbiamo ancora 308 punti di margine che potremmo utilizzare per includere un elemento aggiuntivo che aumenti il valore complessivo della nostra selezione.'\}} \\
\\ 
\texttt{PROPOSTA: \{'B11', 'B26', 'A56', 'A30', 'C47', 'A75', 'A15', 'B91', 'C26'\}} \\
            }
        }
    }
    & & \\ \\

    \theutterance \stepcounter{utterance}  
    & & \multicolumn{4}{p{0.6\linewidth}}{
        \cellcolor[rgb]{0.9,0.9,0.9}{
            \makecell[{{p{\linewidth}}}]{
                \texttt{\tiny{[GM$\rangle$P2]}}
                \texttt{ARGOMENTAZIONE: \{'La tua proposta è solida e rispetta il limite con uno sforzo totale di 3081. Vedo che hai sostituito A19 con B91, il che è una scelta più efficiente. Tuttavia, abbiamo ancora 308 punti di margine che potremmo utilizzare per includere un elemento aggiuntivo che aumenti il valore complessivo della nostra selezione.'\}} \\
\\ 
\texttt{PROPOSTA: \{'B11', 'B26', 'A56', 'A30', 'C47', 'A75', 'A15', 'B91', 'C26'\}} \\
            }
        }
    }
    & \\ \\

    \theutterance \stepcounter{utterance}  
    & & & \multicolumn{4}{p{0.6\linewidth}}{
        \cellcolor[rgb]{0.9,0.9,0.9}{
            \makecell[{{p{\linewidth}}}]{
                \texttt{\tiny{[GM$\langle$P2]}}
                \texttt{ACCORDO: \{'B11', 'B26', 'A56', 'A30', 'C47', 'A75', 'A15', 'B91', 'C26'\}} \\
            }
        }
    }
     \\ \\

    \theutterance \stepcounter{utterance}  
    & & \multicolumn{4}{p{0.6\linewidth}}{
        \cellcolor[rgb]{0.9,0.9,0.9}{
            \makecell[{{p{\linewidth}}}]{
                \texttt{\tiny{[GM$\rangle$P2]}}
                \texttt{La tua risposta non conteneva un argomento. Devi includere il formato strutturato ARGOMENTO: \{...\} da qualche parte nella tua risposta. Riprova.} \\
            }
        }
    }
    & \\ \\

    \theutterance \stepcounter{utterance}  
    & & & \multicolumn{4}{p{0.6\linewidth}}{
        \cellcolor[rgb]{0.9,0.9,0.9}{
            \makecell[{{p{\linewidth}}}]{
                \texttt{\tiny{[GM$\langle$P2]}}
                \texttt{ARGOMENTAZIONE: \{'Perfetto! La tua proposta aggiunge C26 sfruttando efficacemente il margine rimanente. Il set finale con sforzo totale di 3234 rispetta comodamente il limite e rappresenta una selezione ben bilanciata che combina oggetti ad alto valore con elementi efficienti. Questa combinazione ottimizza il nostro utilizzo del budget disponibile.'\}} \\
\\ 
\texttt{ACCORDO: \{'B11', 'B26', 'A56', 'A30', 'C47', 'A75', 'A15', 'B91', 'C26'\}} \\
            }
        }
    }
     \\ \\

    \theutterance \stepcounter{utterance}  
    & & & \multicolumn{2}{p{0.3\linewidth}}{
        \cellcolor[rgb]{0.9,0.9,0.9}{
            \makecell[{{p{\linewidth}}}]{
                \texttt{\tiny{[GM$|$GM]}}
                \texttt{game successful} \\
            }
        }
    }
    & & \\ \\

    \theutterance \stepcounter{utterance}  
    & & & \multicolumn{2}{p{0.3\linewidth}}{
        \cellcolor[rgb]{0.9,0.9,0.9}{
            \makecell[{{p{\linewidth}}}]{
                \texttt{\tiny{[GM$|$GM]}}
                \texttt{end game} \\
            }
        }
    }
    & & \\ \\

\end{supertabular}
}

\end{document}
